% Part A (tech core) of the AI-for-X seminars: limits, mechanisms and security. 4 frames.
% Counterpart of ai_core_nontech_limits (same size, same job, with the mechanism behind each limit).
% Part A is shared: edit here, then recompile every deck that includes it. Trim by commenting frames, not deleting.

%%%%%%%%%%%%%%%%%%%%%%%%%%%%%%%%%%%%%%%%%%%%%%%%%%%%%%%%%%
\begin{frame}[fragile]\frametitle{Why Models Hallucinate}
\begin{itemize}
\item The training goal is a plausible next token, not a true statement. Nothing inside the model checks a fact against a source.
\item Rare facts, exact numbers, quotations, citations and code for obscure libraries are the weak spots.
\item Fluent output carries no signal of correctness, so confidence in the tone means nothing.
\item Mitigations reduce the problem but do not remove it: retrieval (RAG) so the answer comes from supplied text, tools that compute or look up, asking for sources, and automated checks such as tests or schema validation.
\end{itemize}
\end{frame}

%%%%%%%%%%%%%%%%%%%%%%%%%%%%%%%%%%%%%%%%%%%%%%%%%%%%%%%%%%
\begin{frame}[fragile]\frametitle{Tokens, Context Window and Memory}
\begin{itemize}
\item Text is split into \textbf{tokens}, pieces of words. Models read and write tokens, and cost and speed scale with the token count.
\item The model sees only its \textbf{context window}: the prompt, the conversation so far and any documents you add. Anything outside it does not exist for the model.
\item There is no memory between separate chats unless the product stores and re-sends it.
\item Very long contexts can still be used unevenly: details in the middle of a long input are more easily missed. Keep inputs focused.
\end{itemize}
\end{frame}

%%%%%%%%%%%%%%%%%%%%%%%%%%%%%%%%%%%%%%%%%%%%%%%%%%%%%%%%%%
\begin{frame}[fragile]\frametitle{Cutoff, Randomness and Bias}
\begin{itemize}
\item \textbf{Training cutoff:} the weights hold knowledge only up to a date. Newer facts need search or retrieval.
\item \textbf{Sampling:} output is drawn from a probability distribution. A higher temperature gives more varied text, and the same prompt can give different answers. Test prompts on several runs.
\item \textbf{Bias:} training text carries its own biases and gaps, and the model reproduces them.
\item \textbf{Evaluation:} judge a prompt or workflow on a set of examples, not on one good answer.
\end{itemize}
\end{frame}

%%%%%%%%%%%%%%%%%%%%%%%%%%%%%%%%%%%%%%%%%%%%%%%%%%%%%%%%%%
\begin{frame}[fragile]\frametitle{Data Handling and Security}
\begin{itemize}
\item Everything you send goes to the provider. Check its retention and training-use terms for your plan, and prefer the organisation's approved tier or a local model for sensitive material.
\item \textbf{Never paste} secrets, API keys, credentials, personal data or proprietary code into a public tool.
\item \textbf{Prompt injection:} text the model reads (a web page, an email, a document) can contain instructions that the model may follow. Treat retrieved content as untrusted input.
\item Review generated code before running it, and do not give a model tools or permissions you would not give an unreviewed script.
\end{itemize}
\end{frame}
